\section{Benchmark Protocol}
\label{sec:protocol}

The benchmark determines whether a new crack observation matches a recorded defect. Annotated identities stand in for permanent records, so scores indicate association accuracy, not ticket-level duplicate reduction.

\subsection{Task formulation and candidate records}

Every query and its gallery come from one wall, and an identity links all observations of the same crack (Section~\ref{sec:annotation}). Cross-wall search is disallowed; the wall identity is supplied to every method.
A query is a crack instance in one photograph; its gallery contains instances from the other photographs of the same wall.
A query is known if its gallery contains another instance of the same identity, and unknown otherwise (a genuinely new defect).
In closed-set retrieval, the true match is always in the gallery, so candidates are only ranked.
Open-set identification mirrors ticket association: link the query to the best existing ticket or reject it and open a new one.
Method scores drive ranking, thresholded decisions, and identification with rejection.
Pairs are the unit for same-defect decisions; queries are the unit for retrieval and for open-set decisions; image pairs are the unit for registration and assignment.
The Hungarian one-to-one constraint is applied only within each image pair.

\subsection{Test conditions, decisions, and metrics}

The primary test uses a walk-around of each wall, during which crack morphology does not change. The only variations are viewpoint, distance, scale, lighting, and visibility; the galleries also contain same-wall distractors, and some queries are unknown.
The validation and test walls are entirely separate.
All methods are applied off the shelf to the shared masks (Section~\ref{sec:method}), without fine-tuning on \textsc{CrackID} (Section~\ref{sec:setup}).
Thresholds are calibrated on the validation walls and applied unchanged to the test walls; only within-visit correspondence is considered.
A separate synthetic probe isolates growth (condition A: one crack tip extended) from photometric change (condition B: blur, noise, and contrast changes with unchanged structure), comparing OSNet with context (OSNet+ctx, Section~\ref{sec:setup}) against \hybrid{} (Chamfer tolerance $\tau=8.0$~px). Unlike \ours{} (Section~\ref{sec:method}), which performs registration-gated Chamfer comparison only, \hybrid{} blends alignment-based comparison with a crack-skeleton score when homography is available and falls back to skeleton matching otherwise. The probe is preliminary.

\textit{Decisions.} Pairs scoring above the validation-calibrated threshold are accepted. For open-set identification, only the highest-ranked entry is accepted, and only above the threshold; otherwise the query is rejected as a new ticket. Within each image pair, a Hungarian one-to-one assignment is computed and matches below the threshold are discarded. Unscorable pairs (for example, unregistered) are retained: they are predicted negative for $F_1$ and ranked last; a registration failure is absence of evidence, not evidence of a new defect.

\textit{Metrics.} Closed-set ranking uses Rank-1 and mean average precision (mAP); thresholded decisions use pairwise $F_1$ at the oracle threshold ($F_1^{*}$; the test-set $F_1$ maximizer) and at the validation-calibrated threshold ($F_1^{@v}$); assignment uses $F_1^{a} = 2\text{TP}/(2\text{TP}+\text{FP}+\text{FN})$, where TP, FP, and FN denote true positives, false positives, and false negatives. Open-set identification uses the detection and identification rate (DIR) at a false-alarm rate (FAR) of 0.1 (the fraction of known queries correctly linked when at most 10\% of unknown queries are wrongly linked), plus coverage (the fraction of valid query--gallery pairs scored).
Reference levels are random ranking and the constant baseline, which labels every pair as positive and achieves $F_1 = 2p/(1+p)$, where $p$ is positive prevalence. This yields $F_1 = 0.305$ on the full test set ($p = 0.180$), $0.544$ on the geometric subset ($p = 0.373$; $48\%$ of valid pairs), and $0.452$ on validation ($p = 0.292$). Test-set reference levels are in Table~\ref{tab:main}.
Because only 11 walls and 19 matchable identities are available, we report point estimates only (Section~\ref{sec:limitations}).
